\section{Recovery of integral cycles from period jumps}
\label{sec:integral-recovery}

We give the finite-accuracy argument independently of the construction
of the observations. Let $\Lambda=H_3(X_p;\ZZ)$ and let
$\gamma=(\gamma_1,\ldots,\gamma_{134})$ be an integral marking.
Its intersection matrix $J$ is alternating and unimodular. Write
$\delta_i=\gamma d_i$ for the oriented local vanishing spheres and
set
\begin{equation}\label{eq:unknown-map}
 D=(d_1\ \cdots\ d_8):\ZZ^8\longrightarrow\Lambda,
 \qquad n_i=Jd_i,\qquad t_i=\int_{\delta_i}\Psi_p=p_i.
\end{equation}
The two block relations are known; the columns of $D$ are the
unknowns. Their image is primitive by
Proposition~\ref{prop:primitive-vanishing-lattice}; the recovery
procedure below does not require that fact as an input.

In the positive Picard--Lefschetz convention in complex dimension three,
\cite[Expos\'e~XIV, Theorem~3.2.11]{Deligne}, or directly in homology
\cite[equation~(8) in the arXiv-v2 text]{LPPV},
transport around a positive individual node meridian gives
\begin{equation}\label{eq:PL-recovery}
 T_i\gamma_a=\gamma_a+\langle\gamma_a,\delta_i\rangle\delta_i,
 \qquad \Delta_{ai}=(n_i)_a t_i,
 \qquad T_i=I-d_id_i^TJ.
\end{equation}
Here $\Delta_{ai}$ is the endpoint period minus the initial period
of a transported geometric representative of $\gamma_a$. Its forward
evaluation does not require the homology coordinates of that endpoint.
The inversion of \eqref{eq:PL-recovery} is the classical rank-one
Picard--Lefschetz argument; for recovery of homological monodromy from
continued periods see \cite[\S3.5.2, equation~(38), arXiv-v2 text]{LPPV}.
The issue addressed here is its use with
absolute error in both the numerator and the nonzero denominator.

\subsection{A strict stopping criterion}

\begin{lemma}\label{lem:integer-separation}
Suppose $\Delta=nt$ with $n\in\ZZ$ and $t\ne0$. Let
$\widehat\Delta,\widehat t$ satisfy
\[
 |\widehat\Delta-\Delta|\le e_\Delta,
 \qquad |\widehat t-t|\le e.
\]
If $|\widehat t|>e$, define
\begin{equation}\label{eq:observed-height}
 H_{\rm obs}=\frac{|\widehat\Delta|+e_\Delta}
                        {|\widehat t|-e},
 \qquad
 r_{\rm obs}=\frac{e_\Delta+H_{\rm obs}e}{|\widehat t|}.
\end{equation}
If $r_{\rm obs}<1/2$, then $n$ is the unique nearest integer to
$\operatorname{Re}(\widehat\Delta/\widehat t)$.
For any sequence of valid enclosures with $e,e_\Delta\to0$,
the two strict inequalities hold eventually.
\end{lemma}

\begin{proof}
The triangle inequality gives
$|n|=|\Delta|/|t|\le H_{\rm obs}$. Hence
\[
 \left|\frac{\widehat\Delta}{\widehat t}-n\right|
 =\frac{|\widehat\Delta-\Delta+n(t-\widehat t)|}
           {|\widehat t|}
 \le r_{\rm obs}.
\]
Its real part lies strictly within $1/2$ of the integer $n$.
As the errors tend to zero, $|\widehat t|\to|t|>0$,
$H_{\rm obs}\to|n|$, and $r_{\rm obs}\to0$.
\end{proof}

Apply the lemma to the $134\cdot8=1072$ jumps and eight local
periods. Refine all absolute enclosures and accept an entry only after
the displayed strict inequalities have been verified with interval
bounds. A failing strict test need not decide whether equality holds:
one continues the refinement. At a fixed $p\in B_{\rm ev}^\circ$
the lemma proves finite termination for all the entries. There is
no prior bound on the unknown integer heights in this procedure.

The resulting integer vectors give
\begin{equation}\label{eq:recover-D}
                    d_i=J^{-1}n_i\in\ZZ^{134}.
\end{equation}
Both the error in a local period and the absolute error in a small
period difference are indispensable in this step. In particular, a
relative error bound for each of two large endpoint integrals does
not by itself control their small difference.

\subsection{Image, saturation and monodromy}

For this transition $D$ has rank six and
\[
 \ker D=\ZZ v_f\oplus\ZZ v_g,
 \qquad v_f=\sum_{i\in I_f}e_i,
 \qquad v_g=\sum_{i\in I_g}e_i.
\]
The identities $Dv_f=Dv_g=0$ and $D^TJD=0$ are useful consistency
checks. They do not determine the marked inclusion $D$.
Smith reduction of the recovered matrix gives positive integers
$s_1\mid\cdots\mid s_6$ and unimodular changes of source and
target coordinates. If $L=\im D$ and
$L^{\rm sat}=(L\otimes\QQ)\cap\Lambda$, then
\begin{equation}\label{eq:smith-outputs}
 \Lambda/L\cong\ZZ^{128}\oplus\bigoplus_{j=1}^6\ZZ/s_j\ZZ,
 \qquad [L^{\rm sat}:L]=\prod_{j=1}^6s_j.
\end{equation}
Proposition~\ref{prop:primitive-vanishing-lattice} gives
$s_1=\cdots=s_6=1$ and $L^{\rm sat}=L$ before any period is
computed. Smith reduction independently recovers these known
invariants from the observed columns. The columns also determine
the position of $L$ in the specified geometric marking.

\begin{corollary}\label{cor:vanishing-normal-form}
Label $I_f=\{1,2,3,4\}$ and $I_g=\{5,6,7,8\}$.
There is an integral symplectic basis
$\alpha_1,\ldots,\alpha_{67},\beta_1,\ldots,\beta_{67}$ of $\Lambda$
such that
\[
 \begin{aligned}
 \delta_i&=\alpha_i\quad(1\le i\le3),&
 \delta_4&=-\alpha_1-\alpha_2-\alpha_3,\\
 \delta_{4+i}&=\alpha_{3+i}\quad(1\le i\le3),&
 \delta_8&=-\alpha_4-\alpha_5-\alpha_6.
 \end{aligned}
\]
Thus the inclusion and the individual Picard--Lefschetz
transformations are determined up to a simultaneous integral
symplectic change of basis.
\end{corollary}

\begin{proof}
The known kernel makes
$\delta_1,\delta_2,\delta_3,\delta_5,\delta_6,\delta_7$
a basis of $L$. Their disjoint representatives make $L$ isotropic,
and the proposition makes it primitive. Any prescribed basis
$e_1,\ldots,e_6$ of such a sublattice extends to a symplectic
basis. Indeed primitivity allows the coordinate functional
$e_j\mapsto\delta_{1j}$ to extend integrally to $\Lambda$.
Unimodularity represents it by a vector $f_1$ with
$\langle e_j,f_1\rangle=\delta_{1j}$. The remaining $e_j$ lie
in the integral symplectic complement of the plane spanned by
$e_1,f_1$, and span a primitive isotropic sublattice there.
Induction extends the prescribed basis. The two block relations
give the displayed formulas, and \eqref{eq:PL-recovery} gives
the monodromy assertion.
\end{proof}

If $D_0$ denotes this normal-form matrix, then $D=S^{-1}D_0$
for some integral symplectic $S$. Topology determines $D_0$;
the period observations determine $D$ relative to the independently
constructed marking, or equivalently the class of $S$ modulo
the stabilizer of $D_0$.

The local spheres are globally flat by
Section~\ref{sec:transport}, and the marking at $p$ is transported
from $p_*$. Their coordinate matrix is therefore constant under
that transport: $D(p)=D(p_*)$. Recovery at $p_*$ suffices to
determine this matrix throughout $B_{\rm ev}^{\circ}$.
This does not make the full middle-homology local system trivial.

Set $N_i=T_i-I$. The spheres have disjoint representatives, so
$d_i^TJd_j=0$ and therefore $N_iN_j=0$ for all $i,j$. For every
loop $\ell$ in $B_{\rm ev}^\circ$,
\begin{equation}\label{eq:all-monodromy}
 T_\ell=I+\sum_{i=1}^8w_i(\ell)N_i,
 \qquad
 w_i(\ell)=\frac1{2\pi\sqrt{-1}}\int_\ell\frac{\dd p_i}{p_i}.
\end{equation}
Indeed fill the loop by a disc in the convex ball. Perturb its interior
to meet each hyperplane transversely and to avoid pairwise intersections,
which have real codimension four. Removing small discs around the
finitely many inverse images expresses the boundary monodromy as a
product of conjugates of individual meridians. Each local sphere is
fixed by every meridian, so the conjugates have the same action on
homology. The signed numbers of intersections are the windings $w_i$.
Finally all products of two $N_i$ vanish, which proves
\eqref{eq:all-monodromy}. This argument makes no assumption that
the fundamental group of the hyperplane complement is abelian.

\subsection{Precision near a node}

The stopping criterion is adaptive, but a sufficient precision can
also be stated in terms of the true height
$H=\max_a|(n_i)_a|$. Suppose that each initial period, endpoint
period, and local period is evaluated with absolute error at most
$A2^{-b}$, where $A\ge1$. The jump error is at most $2A2^{-b}$.
The conditions
\begin{equation}\label{eq:height-sufficient}
 e_\Delta\le |t_i|/8,
 \qquad e\le |t_i|/[8(H+1)]
\end{equation}
give the direct ratio error
\[
 \left|\frac{\widehat\Delta_{ai}}{\widehat t_i}-(n_i)_a\right|
 \le\frac{e_\Delta+He}{|t_i|-e}<\frac27.
\]
They also imply the strict criterion of
Lemma~\ref{lem:integer-separation}. To see this directly, write
$x=e/|t_i|\le1/[8(H+1)]$. Then
$H_{\rm obs}\le(H+1/4)/(1-2x)$ and
\[
 r_{\rm obs}\le
 \frac{1/8+x(H+1/4)/(1-2x)}{1-x}
 \le\frac27<\frac12.
\]
A sufficient common output precision is consequently
\begin{equation}\label{eq:digit-count}
       b\ge\left\lceil\log_2
                      \frac{16A(H+1)}{|t_i|}\right\rceil.
\end{equation}
The integer height appears only in this explanatory bound, not among
the inputs to the adaptive stopping rule.

For sufficiently small parameters let
$\kappa_i=2\pi/|\log|\tau_i||$. The bounded residue units give
\[
 |t_i|=|\tau_iU_i|\asymp
                     \exp(-2\pi/\kappa_i).
\]
Thus a sufficient number of binary output digits is
\begin{equation}\label{eq:capacity-digits}
 b\ge \frac{2\pi}{\kappa_i\log2}
       +\log_2A+\log_2(H+1)+O(1).
\end{equation}
For direct recovery at the target fibre, this sufficient absolute
tolerance is exponentially small in inverse capacity; the sufficient
number of digits grows linearly in inverse capacity. Recovering the
constant matrix at $p_*$ avoids this target-dependent requirement.
Bounds for transport and quadrature may deteriorate further.
Equation~\eqref{eq:capacity-digits} is neither a running-time bound
nor a bound on the internal precision of those constructions.

\subsection{Changes of choices}

For a target basis change $\gamma'=\gamma S$ with $S$ integral
unimodular,
\begin{equation}\label{eq:target-covariance}
 \begin{aligned}
 J'&=S^TJS,&d_i'&=S^{-1}d_i,& n_i'&=S^Tn_i,\\
 \Delta_i'&=S^T\Delta_i,&T_i'&=S^{-1}T_iS.
 \end{aligned}
\end{equation}
If both bases have the same standard symplectic matrix then $S$ is
symplectic. Absolute observation errors transform with the corresponding
operator norm of $S^T$.

Reversing the orientation of $\delta_i$ changes $d_i,n_i,t_i$ by
the same sign and leaves $\Delta_i,T_i$ unchanged. More generally,
a signed permutation $P$ of the local source basis changes $D$ to
$DP$ and the source relations to $P^{-1}v_f,P^{-1}v_g$. In
particular, the all-positive block relations must not be retained
after an isolated orientation reversal.

A known meridian multiplicity $m\ne0$ multiplies its jump by $m$;
division by $m$ also divides the absolute error by $|m|$. Without
known winding data this primitive jump is not determined. A
single-valued nowhere-zero scalar change of the relative form
multiplies both the local period and the global period jump by the
same scalar at the basepoint, so their ratio is unchanged. If a
chosen form instead has known scalar monodromy, the endpoint must
first be divided by that scalar before the initial period is
subtracted. The canonical residue used here is single-valued.

\subsection{What the observations distinguish}

The marked inclusion cannot be recovered from local periods and the
degree-two spectral data alone. This already follows on one fixed
geometric fibre by changing its middle-dimensional marking. Choose
a primitive $v\in\Lambda$ with $k=v^TJd_i\ne0$ for one nonzero
column. The integral symplectic transformations
\[
 S_N=I-Nvv^TJ,\qquad S_N^{-1}=I+Nvv^TJ
\]
replace that column by $d_i+Nkv$. These columns are unbounded as
$N$ varies, whereas the underlying fibre, local periods, abstract
intersection form and degree-two observations remain the same.
The represented global marking chains do change. Thus this argument
does not produce indistinguishable inputs when the full geometric
marking and all its jumps are retained.

At the scalar level adjacent numerators $nt$ and $(n+1)t$ are
separated by $|t|$; closed error discs of radius $|t|/2$ can meet.
An uncertain denominator gives a further obstruction: $(n,t)$ and
$(n+1,nt/(n+1))$ have the same numerator. These examples explain the
strict half-integer threshold and the role of height in denominator
errors. They are scalar observation obstructions, not sharp minimax
lower bounds for the full geometric inverse problem.

\subsection{The positive reconstruction theorem}

\begin{theorem}\label{thm:integral-reconstruction}
Fix the algebraic central member and six coefficient directions
selected in Section~\ref{sec:algebraic-family}, and the resulting
evaluated family over $B_{\rm ev}$.
For every named $p\in B_{\rm ev}^{\circ}$ there is a terminating
procedure, relative to its coordinate name, which constructs a
geometric symplectic marking of $H_3(X_p;\ZZ)$ and recovers the
marked homomorphism
\[
 D:\ZZ^8\longrightarrow H_3(X_p;\ZZ),\qquad e_i\longmapsto\delta_i.
\]
It returns the primitive integral image and the individual
monodromy matrices, and independently recovers the topologically
determined saturation index one and quotient $\ZZ^{128}$.
The monodromy matrices determine the entire homology
representation by \eqref{eq:all-monodromy}.

The forward observations are eight local periods and $1072$
complex jumps of marked periods around represented individual
meridians. Their absolute errors are obtained from the geometric
construction and refined until
Lemma~\ref{lem:integer-separation} applies. The two source relations
and the ambient symplectic marking are known; the columns of $D$,
their intersection vectors and the endpoint homology coordinates
are not supplied to the forward calculation. For computable $p$
this is an ordinary algorithm. Its sufficient output precision and
its changes under choices are as stated above.
\end{theorem}

\begin{proof}
Proposition~\ref{prop:algebraic-choice} selects the central member;
the rank and algebraic derivative calculation selects the slice.
Theorem~\ref{thm:evaluated-periods} supplies the known inverse
period domain and an algebraic smooth reference fibre.
Theorem~\ref{thm:smooth-marking} constructs the integral marking
there without identifying any local sphere in it.
Theorem~\ref{thm:accessible-observations} transports that marking
to the target, constructs individual meridians, and supplies the
absolute period enclosures by equation-based transport and
integration of represented chains.

The Picard--Lefschetz identity \eqref{eq:PL-recovery} gives
$\Delta_{ai}=(Jd_i)_a p_i$. Every $p_i$ is nonzero at the target.
Lemma~\ref{lem:integer-separation} therefore terminates after
finitely many refinements for every entry, without a height
prior. Unimodularity gives the exact columns by
\eqref{eq:recover-D}; finite integer Smith operations give the
remaining lattice outputs in \eqref{eq:smith-outputs}.
The flat-sphere argument and vanishing pairwise intersections
give \eqref{eq:all-monodromy}. The covariance and precision
statements have been proved above.
\end{proof}

The theorem establishes a construction and termination result.
No numerical coefficient vector, marking, period matrix or attachment
matrix in the constructed marking has been calculated here. The equations
determine the geometry, while the specified procedure obtains the
previously unsupplied marked local inclusion from scalar observations.
